\label{subsec:e-mechanism}
% moved from the main text 2026-08-02 to fit the 8-page limit

To attribute the gain of the full method to its individual components,
rather than to synthetic-contrast augmentation in general, we build an
additive ablation ladder on top of the no-augmentation anchor (\emph{Base},
\cref{subsec:e-main}): first the unsupervised $k$-means intensity partition
alone (\S\ref{subsec:m-transform}, step~1); then the anatomical label-remap
step (step~4); then the Voronoi spatial sub-parcellation (step~2); and
finally the full signed affine remap of \cref{eq:affine} (step~3), which
together constitute PALETTE in its entirety and is denoted \textbf{Ours} in
this ladder --- identical to the standalone (no-Auglab) configuration used
throughout \cref{subsec:e-main}. Every rung is evaluated without Auglab, so
that the ladder isolates PALETTE's own internal components from the
separate question of whether it composes with a prior augmentation
pipeline, which we address next. The synthetic-data fraction is fixed at
$50\%$ of training iterations throughout, as one of several augmentation
constants (alongside, e.g., the $k$-means class count and the affine
magnitude range of \S\ref{subsec:m-transform}) that are set once and not
tuned per dataset; validation uses real-contrast data only (no synthetic
contrast at validation time).

\begin{table}[t]
\centering
\small
\setlength{\tabcolsep}{4pt}
\begin{tabular}{lcccc}
\toprule
Component & CHAOS & Open-MS & BraTS & ON-Harmony \\
\midrule
Base (no augmentation)                & 36.1 & 25.2 & 25.3 & 27.5 \\
\quad + $k$-means partition           & 84.8 & 30.6 & 40.6 & 67.6 \\
\quad + label remap                   & 87.7 & 34.6 & 43.2 & 68.0 \\
\quad + Voronoi sub-parcellation      & 87.2 & 36.1 & 41.8 & 68.2 \\
\quad + affine remap (\textbf{Ours})  & 87.9 & 40.4 & 47.4 & 63.6 \\
\midrule
\quad + Auglab (deployed)             & 87.4 & 39.5 & 46.4 & 68.6 \\
\bottomrule
\end{tabular}
\caption{Mechanism ablation (top five rows, additive, no Auglab) and
deployment composability check (bottom row). Each rung adds exactly one
PALETTE component on top of the row above; all-modality Dice (\%),
averaged over held-out test contrasts, folds 0--2, train fraction $50\%$.
BraTS and ON-Harmony report their T1n/T1w settings, matching every other
table in this paper. The deployed row matches \cref{tab:suppl-dice}'s
\textbf{Ours} column exactly.}
\label{tab:mechanism}
\end{table}

\paragraph{Composability with Auglab.} PALETTE is deployed, throughout this
paper, as one additional operator inside the Auglab pipeline rather than as
a replacement for it (\cref{subsec:e-setup}); Auglab is a full prior
augmentation method in its own right, not a component of PALETTE, so
stacking it on top of the ladder above is a composability/deployment check,
not a further mechanism-ablation step. \cref{tab:mechanism} keeps the two
separate: the top five rows isolate PALETTE's internal mechanism, evaluated
with no Auglab throughout, while the bottom row reports PALETTE plus Auglab
--- the configuration used everywhere else in this paper --- set apart by a
rule rather than presented as the next rung of the same ladder.

